% =============================================================================
\section{Results and Discussion}
\label{sec:results}
% =============================================================================

\subsection{Evaluation Setup}

We evaluate on three captures chosen to stress different assumptions:
\textbf{Mount Rainier} (alpine meadow, dense vegetation, snow),
\textbf{Paris} (flat urban river scene, architecture), and
\textbf{Shark Fin Cove} (coastal cliffs, open water, sea stacks).

Ground truth 3D does not exist for any of them, so evaluation is necessarily
diagnostic: we instrument the pipeline, measure intermediate quantities against
what the scene demonstrably contains, and localize defects to specific stages.

\begin{figure}[htbp]
  \centering
  \begin{subfigure}{0.5\linewidth}
    \centering
    \figbox{rainier-scene.png}{Mount Rainier: reconstructed scene.}{1.0}
    \caption{Mount Rainier}
  \end{subfigure}

  \vspace{0.9em}
  \begin{subfigure}{0.5\linewidth}
    \centering
    \figbox{paris-scene.png}{Paris: reconstructed scene.}{1.0}
    \caption{Paris}
  \end{subfigure}

  \vspace{0.9em}
  \begin{subfigure}{0.5\linewidth}
    \centering
    \figbox{sharkfin-scene.png}{Shark Fin Cove: reconstructed scene.}{1.0}
    \caption{Shark Fin Cove}
  \end{subfigure}
  \caption{Reconstructed environments viewed from near the capture point.}
  \label{fig:scenes}
\end{figure}

\Cref{fig:scenes} shows three reconstructed environments viewed from near the
capture point.

\subsection{Quantitative Summary}

\Cref{tab:results} summarizes scene composition. The mesh/card/billboard split
indicates how far object reconstruction got on each capture, with the caveat that
the card column is dominated by ground cover rather than by object fallbacks.

\begin{table}[t]
\centering
\small
\caption{Scene composition per capture. \emph{Relief added} is the elevation range
introduced by terrain reconstruction beyond what the height field measured.}
\label{tab:results}
\begin{tabular}{@{}lrrrrl@{}}
\toprule
\textbf{Capture} & \textbf{Meshes} & \textbf{Cards} & \textbf{Billboards} &
\textbf{Relief added} & \textbf{Dominant limitation} \\
\midrule
Mount Rainier   & 128 & 9\,303 & 33 & \SI{35}{\meter} & Conifer reconstruction \\
Paris           & 27  & 4      & 56 & \SI{78}{\meter} & Skyline read as landform \\
Shark Fin Cove  & 0   & 1      & 42 & \SI{81}{\meter} & Detections are rock fragments \\
\bottomrule
\end{tabular}
\end{table}

\subsection{What the Measurements Support}

Two behaviors are worth stating separately from the corrective measures, which
\cref{tab:fixes} quantifies in full.

\paragraph{Terrain geometry on landform scenes.}
Where the horizon genuinely is terrain, reconstruction behaves as intended:
measured cells are preserved exactly by construction, the harmonic solve produces a
continuous surface, and erosion supplies drainage-consistent relief. Mount
Rainier's \SI{35}{\meter} of added relief lies beyond depth range, where the
silhouette is the only evidence available and $98\%$ of it types as terrain.

\paragraph{Ground-cover masking.}
Color vetoing excludes snowfields, which type identically to meadow under the
semantic label set. With the veto ordered after morphological closing, zero grass
instances remain on cells it rejected.

\subsection{Failure Modes and Their Causes}

The more instructive result is where the system fails, because in every case we
localized the failure to an \emph{assumption} rather than a model. A defect can
also surface far from its cause: object size compression originates in depth
calibration (\cref{sec:panorama}) and manifests at placement, eight stages later.
\Cref{sec:remediation} reports what each corrective measure achieved.

\begin{figure}[t]
  \centering
  \figbox{paris-terrain.pdf}{Paris height field before and after reconstruction.}{0.92}
  \caption{The Paris failure. \textbf{Left:} the measured height field, spanning
  \SI{-6.0}{\meter} to \SI{-0.2}{\meter}: correctly flat, and near-featureless
  under a shared shading scale. \textbf{Right:} after ridge-anchored reconstruction,
  spanning \SI{-13.1}{\meter} to \SI{+70.9}{\meter}. Hill \textsf{A}
  (\SI{5966}{\meter\squared}) is the right-bank treeline; hill \textsf{B}
  (\SI{2674}{\meter\squared}) is the cathedral. Both panels share one shading
  scale, so the left panel's flatness is the measurement, not a rendering choice.}
  \label{fig:paris-terrain}
\end{figure}

\paragraph{Skyline read as landform.}
Ridge anchoring assumes the sky boundary is the crest of ground. In Paris it is a
roofline \SI{40}{\meter} away. The measured height field spans
\SI{-6.0}{\meter} to \SI{-0.2}{\meter} (flat, correctly), and reconstruction
produced \SI{-13.1}{\meter} to \SI{+70.9}{\meter}: a \SI{5966}{\meter\squared}
hill on the right-bank treeline and a \SI{2674}{\meter\squared} hill on the
cathedral (\cref{fig:paris-terrain}). The discriminator is what the silhouette is \emph{made of}: the fraction
of silhouette columns typing terrain is $0\%$ for Paris against $87\%$, $98\%$ and
$100\%$ for Iceland, Rainier and Shark Fin. Paris is not merely lowest; not one
column of its horizon is a landform (\cref{fig:silhouette}).

\begin{figure}[t]
  \centering
  \figbox{silhouette-composition.pdf}{Terrain share of the sky boundary per capture.}{0.86}
  \caption{What each capture's horizon is made of. The gate at \SI{25}{\percent}
  separates Paris from every landform capture by a wide margin in both directions;
  it is a threshold on evidence the pipeline already computes, not a tuned constant.}
  \label{fig:silhouette}
\end{figure}

\paragraph{Confident misclassification on texture.}
Shark Fin Cove produced 49 detections at median confidence $0.95$; the two
plausibility vetoes of \cref{tab:fixes} later remove 28 of them, and every instance
that survived failed the reconstruction shape gates. The captions are diagnostic: a rock typed
\emph{person} is captioned ``a piece of pizza with bacon on it''; another typed
\emph{aircraft} is ``a giraffe that is flying in the sky''. On captures with weak
image signal the class is derived from the caption ($59\%$ of Shark Fin crops and
$91\%$ of Rainier's), and the captioner describes unparseable texture fluently
and wrongly. These pass label verification, because a tight cutout fills its own
box and an open-vocabulary detector will localize almost any label inside it.

The consequence for reconstruction is direct: SAM~3D given a flat rock fragment
returns a flat sheet (normalized extents $[1.00, 0.049, 0.425]$), the shape gates
correctly reject it, and every instance falls back to a billboard. Shark Fin's zero
meshes are the gates behaving as intended on inputs that are not objects.

The signal that separates them is the region map, because it is the only opinion in
the chain that does not look at the crop. A per-crop plausibility veto, a class
that cannot occupy the region its box sits in, removes $43\%$ of Shark Fin's
detections with no false positives across the other captures.

\paragraph{Thin vegetation reconstruction.}
Conifers fail single-view reconstruction. Given a good $120\times326$ crop, SAM~3D
returns a mesh with aspect $3.50{:}1$ against detections at $0.12{:}1$, a
$28\times$ disagreement, so 70 of 75 instances fall back to crossed cards. A thin
conifer offers almost no volume cue from one view. Species-aware
approaches~\cite{lee2024tree} are the natural remedy.

\paragraph{Water surface following terrain.}
Water inherited the terrain height under each of its cells, which is correct only
if the ground under a lake is flat. On Shark Fin one \SI{11872}{\meter\squared}
body of sea spanned \SI{73.5}{\meter} in elevation, with $21.7\%$ of its vertices
more than a meter above the median: the ocean draped up the cliffs.

\begin{figure}[t]
  \centering
  \figbox{texel-density.pdf}{Terrain texel density by distance ring.}{0.86}
  \caption{Texture resolution against distance, log scale. The two captures are
  indistinguishable within \SI{25}{\meter}; beyond it Shark Fin degrades far
  faster, because a clifftop camera places most of its surface near the horizon
  where equirectangular compression is worst. Above roughly
  \SI{0.5}{\meter} per texel we take the ground to read as smeared.}
  \label{fig:texel}
\end{figure}

\paragraph{Far-field texture resolution.}
Even with folds repaired, $67.8\%$ of Shark Fin's terrain area exceeds
\SI{0.5}{\meter} per texel against $20.9\%$ for Rainier, because a clifftop camera
places most of its surface near the horizon where equirectangular compression is
worst (\cref{fig:texel}).

\subsection{Corrective Measures and Their Measured Effect}
\label{sec:remediation}

Each failure above was addressed by making its implicit premise testable from
evidence already computed elsewhere in the pipeline, rather than by tuning a
threshold or replacing a model. \Cref{tab:fixes} reports the measured effect of
each, evaluated by replaying the modified stage over saved pipeline state from all
five captures. We report effects on \emph{every} capture, not only the one that
motivated the change, because the risk being managed is regression on scenes that
were already correct.

\begin{table}[t]
\centering
\small
\caption{Corrective measures, the premise each makes testable, and measured effect
across captures. ``FP'' denotes false positives: correct content wrongly removed.}
\label{tab:fixes}
\begin{tabular}{@{}p{3.3cm}p{4.5cm}p{6.2cm}@{}}
\toprule
\textbf{Measure} & \textbf{Premise made testable} & \textbf{Measured effect} \\
\midrule
Silhouette composition gate &
``The sky boundary is the crest of ground'' &
Terrain-typed silhouette fraction: Paris $0\%$; Iceland $87\%$; Rainier $98\%$;
Shark Fin $100\%$. Ridge anchoring suppressed on Paris only; other captures
bit-identical. \\
\addlinespace
Per-crop region plausibility &
``This class can occupy the region its box sits in'' &
Removes 9/109 (Paris), 4/237 (Rainier), 12/49 (Shark Fin), 5/101 (Iceland).
0 FP; Rainier's four are phantom aircraft inside the snowfield. \\
\addlinespace
Caption--scene plausibility &
``The caption describes something this scene can contain'' &
Removes a further 16 (Shark Fin), 16 (Iceland), 9 (Rainier), 2 (Paris). 0 FP.
Requires four conjunctive conditions; each weaker variant deletes real content
(see text). \\
\addlinespace
Water surface levelling &
``The ground beneath water is flat'' &
Shark Fin elevation spread \SI{73.5}{\meter} $\rightarrow$ \SI{0.00}{\meter};
resulting shoreline discontinuity \SI{0.14}{\meter}. Paris (already near-level)
\SI{1.67}{\meter} $\rightarrow$ \SI{0.55}{\meter} across 12 bodies. \\
\addlinespace
Projected-size LOD &
``Distance predicts resolvability'' &
Rainier meshed instances $545 \rightarrow 296$; triangles
\SI{4.47}{M} $\rightarrow$ \SI{2.48}{M}; per-class override table eliminated. \\
\addlinespace
Resolution-based asset selection &
``The best-scoring crop is the best reconstruction input'' &
Source crop area increased $1.0$--$9.2\times$ per group; \num{205} of \num{471}
instances had been rejected as degenerate reconstructions. \\
\addlinespace
Geometric UV-fold repair &
``Adjacent grid cells observed adjacent panorama rows'' &
Collapses \num{5624} triangles (Rainier, $7.0\%$), \num{2169} (Shark Fin),
\num{1479} (Paris). Converges in two passes; face count unchanged. \\
\addlinespace
Occluder-footprint inpainting &
``What foreground removal left behind is ground'' &
Below-horizon pixels showing fabricated fill: $23.9\% \rightarrow 17.1\%$. \\
\bottomrule
\end{tabular}
\end{table}

Two of these deserve comment because the obvious version of each is measurably
wrong.

\paragraph{Caption plausibility requires conjunction.}
Rejecting a detection because its caption mentions something implausible is
appealing and unsafe. On Mount Rainier, 21 flower detections carry captions such as
``a close up of a flower in a vase'', and \num{19} of those \num{21} are real,
placed flowers. The captioner identifies the \emph{subject} correctly and
confabulates the \emph{setting}. Conversely, rejecting on scene-tag support alone
is also unsafe: Rainier's tags name no tree, and it has \num{67} real ones. Only
the conjunction of four conditions (no scene-tag support, implausible caption
vocabulary, caption not naming the class, and class not vegetation) produced
zero false positives. Notably, the caption--class agreement test is unusable as a
veto (its synonym gaps delete real objects) but safe as an \emph{exemption}, because
a gap then merely retains a spurious detection rather than removing a real one.

\paragraph{Robust estimation must respect sampling density.}
Water levelling initially estimated each body's surface from its mesh vertices.
Because the mesh is Poisson-disc sampled with a near-camera density bias,
\num{4309} of Shark Fin's \num{11366} water vertices lie within \SI{20}{\meter} of
the camera, and the resulting estimate described the cove at the viewer's feet
rather than the sea: \SI{-4.54}{\meter} against a shoreline at
\SI{-0.99}{\meter}. Estimating over the uniform grid instead gives
\SI{-1.13}{\meter} and a \SI{0.14}{\meter} shoreline step. The estimator was
correct; the sample weighting was not.
